\chapter{Conclusions and Future Work}
\label{ch:conclusions}


This chapter answers the three parts of the question posed in Chapter~\ref{ch:introduction},
separating what was established from what was argued, and sets out the future work that
follows from the limitations of Chapter~\ref{ch:evaluation}.

\section{Conclusions}

This project asked where strategic informational advantage sits in a prediction market whose
traded object is a belief and whose resolver is a model, what it is worth, and what it
should be called. It answers the three parts unevenly, and the unevenness is the honest
result.

On \emph{where}, the Sapience half of the answer is established and the Delphi half only
partly. A Sapience request is broadcast in full to every client connected to the relayer,
with no allowlist and no filtering, carrying the direction of every leg, the position size
and the requester's address, and a forecaster who wants a fill has no way to send less. That
comes from the protocol's own documentation and implementation, and the protocol does not
conceal it; the Sapience trust model addresses custody, forgery and authority and simply
does not treat observation as a harm. Delphi adds no pre-trade disclosure of its own,
exposing intent only through a trade's effect on a pooled price, though its on-chain
transactions remain subject to the ordinary transaction-ordering surface; at resolution it
relocates the question to who specifies what the market meant. That last part rests on a
design Gensyn documented in May 2026 whose current status their own documentation does not
establish, as Section~\ref{sec:limitations}(h) records.

On \emph{what it is worth}, the answer is measured but bounded. Against a matched benchmark
holding competition, signal quality, dependence and risk identical, letting responders
condition on the request removes about seventy-three per cent of an informed forecaster's
surplus at the baseline configuration. Direction alone accounts for roughly ninety-five per
cent of that, and the share lost rises with the forecaster's accuracy to about ninety-six
per cent at the highest accuracy modelled. Those are properties of a simulator, not of a
venue. The model contains no uninformed order flow, so its responders cannot be solvent in
any configuration and it has no participation constraint; the magnitudes are therefore upper
bounds on the forecaster's position rather than estimates of anyone's.

On \emph{whether it is AI-MEV}, the honest answer is that this project does not settle it,
and the reasons are worth more than a verdict would have been. The second limb of the
definition, that the advantage exceeds ordinary compensation, is satisfied. The first, that
the access is role-specific, is not, because both of the conditions testing it turn on whether capital-gated access to a public stream counts as
protocol-specific, which is a question about the definition rather than about Sapience, and
one this report inherits from the MEV literature rather than introducing. The condition that
would make the case deepest, accumulated advantage from a participant's history, is argued
from protocol structure and never measured. And the exclusion for risk compensation is set
to zero by the model's assumption of risk-neutral responders rather than tested against
anything.

What does survive, and does not depend on the verdict, is a design critique. The mechanism
compels complete, address-linked disclosure and offers nothing cheaper. Measured against the
matched benchmark, the entire measured transfer accrues to the responders who quote on it,
rather than to every party that merely watches the stream. And the cost rises with the quality of the information disclosed. A
prediction market whose disclosure cost scales with forecaster skill is taxing the input it
exists to attract, whatever the transfer is called.

The report's other contribution needs no experiment and is the more durable of the two.
Semantic validity, reproducibility and execution attestation are three properties, not three
names for one. Reproducibility and attestation are not substitutes, because frontier models
are non-deterministic and an inference cannot be proved by repeating the call. Neither
touches whether the committed specification was a fair rendering of the question, and
nothing in either protocol's stack does. Delphi's own history supports this: a designer with
every incentive to close the whole problem removed the settling party's execution
discretion, and published nothing to indicate that specification discretion went with it.
That is an argument from silence, and the silence has a competing explanation, since Gensyn
document no settlement mechanism at all; it is offered as support and not as proof.
Verification is
a statement about a computation; the residual is a statement about meaning, and it does not
yield to more verification.

The two halves meet at a single observation. Both protocols have built serious machinery
against dishonesty, and what survives that machinery in both cases is not dishonesty. At
execution it is a disclosure the mechanism compels and does not price. At resolution it is
an ambiguity the mechanism cannot adjudicate. Neither is anyone's misconduct, which is
precisely why neither is fixed by the mechanisms designed to detect misconduct.

\section{Future work}

Most of the items below follow directly from the lettered limitations in
Section~\ref{sec:limitations}; the last two are refinements the analysis suggested rather
than gaps it left, and are labelled as such.

\paragraph{Responders who price risk and can decline.} The deepest limitation is that this
model cannot separate information rent from risk compensation, because its responders are
risk-neutral and quote their own posteriors without shading. A model in which responders
set a protective spread, hold inventory and may refuse to quote would make that separation measurable rather
than assumed, and would also produce a market whose participants do not lose money. This is
the single change that would most improve the empirical half of the project.

\paragraph{Accumulation across requests.} The accumulation limb is the report's weakest and
the one a critic should attack first. Giving responders memory across events, so that an
observer builds a posterior over a forecaster's skill from a reused address, would test
whether the advantage this report argues for exists at all, and whether it survives a
forecaster who rotates addresses. The concentration parameter specified in the project's
design was not implemented and should be.

\paragraph{The forecaster's own strategic choices.} Limitations (d), (e) and (f) each
remove a decision a real forecaster makes: whether to submit a request at all when the
expected disclosure cost exceeds the expected gain, how much to stake, and whether to split
a position across requests. All three are cheap to add to the present model, and the first
is plausibly the most important strategic response available to a real participant. The
size-to-confidence correspondence on which the $B$ sweep's mapping to a real venue rests,
limitation (e), is also testable against data if any venue publishes it, and it should be
tested before that mapping is relied on.

\paragraph{Multi-leg disclosure.} Chapter~\ref{ch:protocols} argues that a combination
reveals joint rather than marginal belief, which is the strongest form of the disclosure
argument, and no experiment supports it. Legs with a controlled redundancy parameter would
test whether disclosure scales with the number of legs or only with their informational
independence.

\paragraph{A calibrated comparison against a pari-mutuel arm.} Running the same informed
forecaster through a dynamic pari-mutuel mechanism would let the two protocols be compared
directly rather than analysed side by side. It was scoped out, as Section~\ref{sec:sweep}
records, because the calibration is a research question in its own right: comparing an auction with a pool
requires matching capital at risk and transaction cost, and a comparison done carelessly
would produce a headline about which protocol is better that the evidence could not support.

\paragraph{Specification sensitivity.} The proposition that no verification mechanism
guarantees semantic validity is argued from protocol analysis. It could be measured: holding execution fixed and
varying a market's settlement specification in ways a reasonable creator might, one could
report the rate at which the resolved outcome changes. The methodological care required is
substantial, since a settled market yields a recorded protocol outcome rather than a ground
truth, and for the ambiguous questions that are actually interesting the outcome is part of
what is under study. This was scoped out for time, as Section~\ref{sec:sweep} records, and
remains the natural next experiment.

\paragraph{Source integrity as a fourth property.} A refinement rather than a gap. This report treats semantic validity as
one property, but whether the evidence a judge read was accurate and unaltered is arguably a
fourth, distinct from all three: a perfectly worded question can be settled against corrupted
evidence, and neither reproducibility nor attestation speaks to it. Gensyn's own statement
that receipts do not execute tools and do not prove an external tool returned accurate data
\cite{gensyn2026receipts} is direct evidence for the distinction. It is not adopted here
because populating the category properly would require analysis of data-source provenance in
AI settlement that this project has not done, and naming a category one cannot then fill is
worse than naming the question.

\paragraph{Conditional recall as a mitigation.} Also a refinement. Evaluating a request inside a trusted
execution environment and requiring it to be discarded when no trade results would target
the accumulation gap directly. It is the mitigation best matched to the limb this report
could not measure, and it is testable: an enclave-based responder could be compared against
an unconstrained one under the accumulation model described above.

\section{Closing remarks}

The project set out to find where value is extracted in agentic prediction markets. In the
half it could measure it found a transfer it can bound and cannot classify. Three things
obstruct the classification and only one of them is external: the risk exclusion is zeroed
by this model's assumptions, accumulation was never measured, and the question of whether
capital-gated access to a public stream is protocol-specific is inherited from the MEV
literature. Two of the three are mine. Saying so
is more useful than a verdict the evidence would not carry. The claim that survives is
smaller and better supported: a venue that leaves a participant no cheaper way to trade than
to say everything has made a design choice, that choice has a price, and the price is
highest for the participants the venue most wants.

What I would defend furthest is the part that required no experiment. As settlement moves
from people to models, the instinct is to ask whether the computation can be verified, and
increasingly it can be. The question that then remains is whether the computation was the
right one to perform, and no receipt, transcript or bond answers it. A protocol can remove
every class of privilege it is able to name and still leave the party who wrote the question
holding the only power that finally decides the outcome.
